\section{Оптимальное расписание пакетов}
\label{sec:batch}

\subsection{Точная стоимость для нулевого учителя}

Пусть $Q=\Bern(0)$, а априорное распределение вероятности единицы равно
$\Beta(a,b)$, где $a,b>0$. После $n$ нулей апостериорное распределение имеет вид
$\Beta(a,b+n)$. Одношаговый байесовский предиктор и его риск равны
\begin{equation}
 P_n(1)=\frac{a}{n+a+b},\qquad
 P_n(0)=\frac{n+b}{n+a+b},\qquad
 K_n=\ln\left(1+\frac{a}{n+b}\right).
 \label{eq:beta-risk}
\end{equation}
Следовательно, минимальный объём для риска не выше $\varepsilon$ равен
\begin{equation}
 N_\varepsilon=\max\left\{0,
  \left\lceil\frac{a}{e^\varepsilon-1}-b\right\rceil\right\}
   \sim\frac{a}{\varepsilon}.
 \label{eq:target-N}
\end{equation}

Учитель выбирает только полностью нулевой пакет. В PFR минимум среди таких
кандидатов достигается на первом из них: их отношения плотностей одинаковы,
а времена прихода строго возрастают. При замороженном предикторе вероятность
успеха одного кандидата равна $P_n(0)^B=e^{-BK_n}$. Поэтому
\begin{equation}
 J\sim\Geom(e^{-h}),\quad h=BK_n,\qquad
 c_{\rm fr}(n,B)=L(BK_n),\quad
 L(h)=\E_{J\sim\Geom(e^{-h})}\ell_\delta(J).
 \label{eq:exact-transition}
\end{equation}
Здесь геометрическое распределение имеет носитель $\{1,2,\dots\}$, так что
$\E J=e^h$. В данной специальной модели индекс можно найти простым поиском
первого успеха; общий PFR-поиск сложнее.

\begin{lemma}[Средняя длина геометрического индекса]
\label{lem:length}
Функция $L$ непрерывна на $[0,\infty)$, не убывает, $L(0)=1$ и является
непрерывно дифференцируемой на $(0,\infty)$. Для $h\ge0$
\begin{equation}
 L(h)=1+\sum_{j=1}^{\infty}d_j(1-e^{-h})^{2^j-1},\qquad
 d_j=\begin{cases}3,&j+1\text{ --- степень двойки},\\1,&\text{иначе}.
 \end{cases}
 \label{eq:dyadic-series}
\end{equation}
Кроме того,
\begin{equation}
 L(h)\ge\max\{1,h/\ln2\},\qquad
 L(h)\le\frac{h}{\ln2}
      +2\log_2\left(1+\frac{h}{\ln2}\right)+1.
 \label{eq:L-bounds}
\end{equation}
\end{lemma}
\begin{proof}
Длина~\eqref{eq:elias} меняется только при переходе индекса через $2^j$.
Приращение равно $d_j$. Формула ожидания через вероятности превышения порогов
даёт~\eqref{eq:dyadic-series}. На каждом отрезке $0\le h\le H<\infty$ ряд
равномерно сходится, поскольку его члены ограничены величинами
$3(1-e^{-H})^{2^j-1}$. На компактных подотрезках $(0,\infty)$ так же сходится
ряд производных. Монотонность следует из монотонности каждого слагаемого.

Нижняя граница единицей очевидна. Для $h>0$ вероятность любого значения
геометрического индекса не выше $e^{-h}$, поэтому его энтропия не меньше
$h/\ln2$. Средняя длина префиксного кода не меньше энтропии. Для верхней
границы используем
$\ell_\delta(j)\le\log_2j+2\log_2(1+\log_2j)+1$ и дважды неравенство Йенсена:
$\E\log_2J\le\log_2\E J=h/\ln2$.
\end{proof}

Обозначим через $V_{a,b}(N)$ минимальную ожидаемую сумму длин отдельных индексов
среди публичных расписаний, заканчивающихся ровно на $N$. Через
$C_{\rm fr}(\varepsilon)$ обозначим минимум при условии
$K_N\le\varepsilon$. Расписания конечны; при $N=0$ сообщения не требуются.

\begin{theorem}[Ведущий коэффициент и ограниченный остаток]
\label{thm:optimal}
Для любых $a,b>0$ положительный конечный минимум
\begin{equation}
 A(a)=\min_{r>1}\frac{L(a(r-1))}{\ln r}
 \label{eq:A}
\end{equation}
достигается. При $N\to\infty$ и $\varepsilon\downarrow0$ справедливы
\begin{align}
 V_{a,b}(N)&=A(a)\ln(N+b)+O_{a,b}(1),
 \label{eq:finite-asymp}\\
 C_{\rm fr}(\varepsilon)&=A(a)\ln(1/\varepsilon)+O_{a,b}(1).
 \label{eq:eps-asymp}
\end{align}
Для каждого $N$ имеется явная нижняя оценка
\begin{equation}
 V_{a,b}(N)\ge
 A(a)\ln\frac{N+b+a}{b+a}.
 \label{eq:finite-lower}
\end{equation}
Геометрический рост $n+b$ с любым минимизатором $r_*$
в~\eqref{eq:A}, с целочисленным округлением и усечением последнего пакета,
достигает~\eqref{eq:finite-asymp}.
\end{theorem}
\begin{proof}
При $r\downarrow1$ отношение в~\eqref{eq:A} стремится к бесконечности, так как
$L\ge1$. При $r\to\infty$ оно не меньше
$a(r-1)/(\ln2\ln r)$ и также стремится к бесконечности. Непрерывность $L$
обеспечивает достижение положительного конечного минимума.

Из $\ln(1+x)\ge x/(1+x)$ и монотонности $L$ для любого перехода следует
\[
 L(BK_n)\ge L\left(\frac{aB}{n+b+a}\right)
 \ge A(a)\ln\left(1+\frac{B}{n+b+a}\right).
\]
Сумма логарифмов по расписанию телескопируется, что
даёт~\eqref{eq:finite-lower}.

Для верхней оценки зафиксируем минимизатор $r_*>1$ и положим на бесконечном
расписании
\[
 B(n)=\max\{1,\lfloor(r_*-1)(n+b)\rfloor\}.
\]
Пишем $x=n+b$. После конечного числа начальных шагов
\[
 B=(r_*-1)x+O(1),\quad
 \ln\frac{x+B}{x}=\ln r_*+O(x^{-1}),\quad
 B\ln(1+a/x)=a(r_*-1)+O(x^{-1}).
\]
Величина $x$ растёт геометрически, поэтому сумма $x^{-1}$ по шагам конечна.
Лемма~\ref{lem:length} даёт локальную липшицевость $L$ в окрестности
$a(r_*-1)>0$. После $m$ полных пакетов суммарная стоимость равна
$mL(a(r_*-1))+O(1)$, а прирост $\ln x$ --- $m\ln r_*+O(1)$.
Перед первым превышением $N$ укорачиваем последний пакет. Его стоимость
ограничена константой, как и отличие $\ln x$ от логарифма конечной точки
полного пакета. Получаем верхнюю оценку в~\eqref{eq:finite-asymp}.

Любое расписание до $N'>N$ можно усечь на $N$: последняя стоимость уменьшится
по монотонности $L$. Поэтому $V_{a,b}(N)$ не убывает и
$C_{\rm fr}(\varepsilon)=V_{a,b}(N_\varepsilon)$.
Подстановка~\eqref{eq:target-N} завершает доказательство.
\end{proof}

Параметр $a$ имеет также геометрическую интерпретацию. Если $p=w^k$ при
равномерном $w\in[0,1]$, то $p\sim\Beta(1/k,1)$ и $a=1/k$.
Для нулевого учителя $K(w)=-\ln(1-w^k)\sim w^k$, так что априорная масса
множества $\{K(w)<\varepsilon\}$ имеет порядок $\varepsilon^{1/k}$.
В этой граничной модели показатель априорной массы равен $a$.
Теорема~\ref{thm:optimal} устанавливает его связь с длиной записи через
точную формулу~\eqref{eq:beta-risk} и справедлива для любого вещественного
$a>0$.

\begin{proposition}[Конечный оптимум]
\label{prop:dp}
Минимальная длина до фиксированного $N$ вычисляется рекурсией
\begin{equation}
 V_{a,b}(0)=0,\qquad
 V_{a,b}(j)=\min_{0\le n<j}
 \left\{V_{a,b}(n)+L\left((j-n)\ln\left(1+\frac a{n+b}\right)\right)\right\}.
 \label{eq:dp}
\end{equation}
Она требует $O(N^2)$ вычислений стоимости перехода и $O(N)$ памяти для
значений и предшественников. Хранение всей матрицы стоимостей дополнительно
требует $O(N^2)$ памяти.
\end{proposition}
\begin{proof}
Последний переход оптимального пути в вершину $j$ начинается в некотором
$n<j$. Его предыстория также должна быть оптимальной, иначе её можно заменить
на более дешёвую. Обратное утверждение следует приписыванием последнего
перехода. Получается кратчайший путь в ациклическом графе объёмов выборки.
\end{proof}

\subsection{Ограничение числа кандидатов}

Оценка длины не учитывает стоимость поиска индекса. В нашей модели ожидаемое
число просмотренных кандидатов на шаге в точности равно $e^{BK_n}$.
Потребуем, чтобы оно не превышало публичную константу $M$ на каждом шаге.
Непустое расписание существует тогда и только тогда, когда
$M\ge1+a/b$: уже первый пакет содержит не менее одного наблюдения.

\begin{theorem}[Оптимум при ограничении поиска]
\label{thm:compute}
Пусть $M\ge1+a/b$. Для оптимальной ожидаемой длины среди расписаний с
$\E J_i\le M$ на каждом шаге справедливо
\begin{equation}
 C_{{\rm fr},M}(\varepsilon)
 =A_M(a)\ln(1/\varepsilon)+O_{a,b,M}(1),\qquad
 A_M(a)=\min_{1<r\le1+\ln M/a}
       \frac{L(a(r-1))}{\ln r}.
 \label{eq:compute-opt}
\end{equation}
Коэффициент $A_M(a)$ не возрастает по $M$. Если область минимизации содержит
хотя бы один минимизатор~\eqref{eq:A}, то $A_M(a)=A(a)$.
\end{theorem}
\begin{proof}
Для любого допустимого перехода
\[
 \frac{aB}{n+b+a}\le BK_n\le\ln M.
\]
Следовательно, $r=1+B/(n+b+a)$ принадлежит области в~\eqref{eq:compute-opt}.
Нижняя оценка из доказательства теоремы~\ref{thm:optimal} сохраняется с $A_M$.

Для верхней оценки выбираем минимизатор $r_M$ и $h_*=a(r_M-1)\le\ln M$.
Используем правило
\[
 B(n)=\max\{1,\lfloor h_*/K_n\rfloor\}.
\]
Если округление даёт ноль, пакет из одного наблюдения допустим, поскольку
$K_n\le K_0\le\ln M$. После конечного начального участка
$BK_n=h_*+O((n+b)^{-1})$ и
$B/(n+b)=h_*/a+O((n+b)^{-1})$.
Поэтому стоимость и логарифмический прирост имеют те же суммируемые остатки,
что и в предыдущем доказательстве. Последний пакет можно укоротить без
нарушения ограничения. Монотонность $A_M$ следует из вложенности областей.
\end{proof}

При таком расписании число сообщений имеет порядок $\ln(1/\varepsilon)$.
Такой же порядок имеет ожидаемое суммарное число просмотренных пакетов
из потока кандидатов.
Если генерация одного пакета из $B$ наблюдений стоит $O(B)$, суммарная
ожидаемая стоимость их генерации равна $O(N_\varepsilon)=O(1/\varepsilon)$:
$\sum_i B_i\E J_i\le MN_\varepsilon$. Это утверждение относится к генерации
кандидатов; стоимость обучения модели учитывается отдельно.

\subsection{Совместный предиктив вместо замороженного}

У байесовского ученика имеется другой естественный способ генерировать пакет:
выбрать параметр из апостериорного распределения и затем получить $B$
условно независимых меток.
Такое распределение совпадает с последовательным предсказанием с
внутрипакетным обновлением. Оно отличается от $P_n^{\otimes B}$.
Для полностью нулевого пакета его вероятность равна
\begin{equation}
 P_n^{\rm joint}(0^B)=\frac{\mathrm B(a,b+n+B)}{\mathrm B(a,b+n)},\qquad
 h_{\rm joint}(n,B)=\ln\frac{\mathrm B(a,b+n)}{\mathrm B(a,b+n+B)},
 \label{eq:joint}
\end{equation}
где $\mathrm B$ --- бета-функция. Индекс первого нулевого кандидата по-прежнему
геометрический, а его стоимость равна $L(h_{\rm joint}(n,B))$.
Конечное апостериорное распределение после $N$ нулей то же, что у замороженного
протокола.

\begin{theorem}[Зависимость ведущей константы от предиктива]
\label{thm:joint}
Обозначим через $C_{\rm joint}(\varepsilon)$ оптимальную ожидаемую длину
записи с предложением~\eqref{eq:joint}, теми же априорным распределением,
кодом индекса и условием
конечного риска. При $\varepsilon\downarrow0$
\begin{equation}
 C_{\rm joint}(\varepsilon)=\frac a{\ln2}\ln(1/\varepsilon)
                  +O(\ln\ln(1/\varepsilon)),
 \qquad \frac a{\ln2}<A(a).
 \label{eq:joint-asymp}
\end{equation}
Один пакет из $N_\varepsilon$ наблюдений уже достигает этой ведущей константы.
\end{theorem}
\begin{proof}
Для любого расписания приращения~\eqref{eq:joint} телескопируются:
\[
 \sum_i h_{\rm joint}(n_i,B_i)=
 H_N:=\ln\frac{\mathrm B(a,b)}{\mathrm B(a,b+N)}
      =a\ln N+O_{a,b}(1).
\]
Последнее равенство следует из
$\Gamma(N+a+b)/\Gamma(N+b)\sim N^a$.
По~\eqref{eq:L-bounds} стоимость любого расписания не меньше $H_N/\ln2$.
Для одного пакета её значение равно $L(H_N)$ и не превосходит
$H_N/\ln2+2\log_2(1+H_N/\ln2)+1$.
Поскольку $H_N$ возрастает, условие $N\ge N_\varepsilon$ и
\eqref{eq:target-N} дают требуемую асимптотику.

Для замороженного протокола минимум достигается в конечной точке $r_*>1$.
Используя $L(h)\ge h/\ln2$ и $r_*-1>\ln r_*$, получаем
\[
 A(a)\ge\frac{a(r_*-1)}{\ln2\ln r_*}>\frac a{\ln2}.
\]
\end{proof}

Выигрыш в длине сопровождается ростом поиска: у одного совместного пакета
$\E J=e^{H_N}=\Theta(N^a)$, тогда как геометрическое расписание замороженных
пакетов удерживает эту величину ограниченной. Теорема~\ref{thm:joint}
не утверждает оптимальность одного пакета по времени.

\begin{remark}[Область нижних оценок]
Нулевой учитель выбран для точного анализа. Если декодеру разрешено пропустить
индексы и просто получить $N$ нулей, он воспроизведёт апостериорное распределение без сообщения
при публичном $N$. Поэтому положительные константы выше характеризуют класс
протоколов с \emph{обязательной отдельной передачей индекса каждого пакета}.
Они не являются нижними границами для всех способов передачи этой модели.
\end{remark}

\subsection{Случайный учитель и ожидаемая верхняя оценка}

Для $Q=\Bern(q)$, $0<q<1$, состояние после $n$ выбранных примеров определяется
числом единиц $S_n\sim\operatorname{Bin}(n,q)$. Условно на $S_n=s$
вероятность единицы ученика равна $p=(s+a)/(n+a+b)$.
Для фиксированного публичного расписания стоимость перехода усредняется по
этому биномиальному распределению; размеры пакетов фиксируются независимо
от реализованных значений $s$.

Для распределений пакетов $P,Q$ положим $R(x)=Q(x)/P(x)$. Используя известный
условный закон PFR-индекса \citep[формулы (34)--(38)]{qiu2026}, получаем
\begin{equation}
 J\mid X=x\sim\Geom(M(x)^{-1}),\qquad
 M(x)=1+\sum_y P(y)(R(x)-R(y))_+.
 \label{eq:conditional-index}
\end{equation}
Отсюда точная ожидаемая стоимость равна
$\E_{S_n}\E_{X\sim Q^{\otimes B}}L(\ln M(X))$.
В бернуллиевской модели внутренняя сумма сводится к числу единиц в пакете,
так что все ожидания являются конечными биномиальными суммами. Рекурсия
\eqref{eq:dp} с этой усреднённой стоимостью находит оптимальное публичное
расписание до фиксированного $N$.

\begin{proposition}[Логарифмический порядок для суррогатной стоимости]
\label{prop:surrogate}
Пусть при любом детерминированном расписании распределение ученика после $n$
примеров одинаково, $K_n\ge0$, и для некоторых $0<c_1\le c_2<\infty$
\[
 \frac{c_1}{n+1}\le\E K_n\le\frac{c_2}{n+1}\quad(n\ge0).
\]
Минимальная сумма \emph{верхних оценок}
$\sum_i\E f(B_iK_{n_i}/\ln2)$ при $\E K_N\le\varepsilon$ имеет порядок
$\Theta(\ln(1/\varepsilon))$.
\end{proposition}
\begin{proof}
Так как $f(d)\ge d$, стоимость не меньше
\[
 \frac{c_1}{\ln2}\sum_i\frac{B_i}{n_i+1}
 \ge\frac{c_1}{\ln2}\ln(N+1).
\]
Условие риска требует $N+1\ge c_1/\varepsilon$. Для верхней оценки выбираем
удвоения $n+1$ и конечный объём порядка $c_2/\varepsilon$.
Тогда $B_i\E K_{n_i}$ ограничены, а число пакетов логарифмическое.
Вогнутость $f$ и неравенство Йенсена ограничивают каждый член сверху
константой.
\end{proof}

Доказательство не даёт нижнюю оценку фактической PFR-длины: из нижней
границы на верхнюю оценку нельзя вывести нижнюю границу на исходную величину.
Для шумного учителя точная асимптотическая константа оптимальной длины здесь
не установлена. На конечном горизонте точная стоимость и
\eqref{eq:universal-bound} могут выбирать разные расписания
(раздел~\ref{sec:validation}).
